\subsection{研究问题}

本文的实验设计围绕五个问题展开。第一，VLM 在单图条件下是否真的具备稳定的局部空间关系理解能力，而不是只会利用语言模板或局部纹理线索。第二，多视角输入是否会把单视角中的脆弱正确率整合为更稳定的场景表征，还是会暴露模型缺乏统一场景信念的问题。第三，模型给出的 QRR/TRR 回答能否被联合解释为一个可实现、可重建、与真实场景对齐的对象级结构。第四，这种结构一旦放入时间维度，在物体运动、遮挡和关系翻转出现后能否继续保持一致。第五，当引入 CoT、visual CoT、工具使用或 test-time search 后，模型提升的究竟是“题目得分”，还是“场景级可重建性”。

\subsection{静态场景协议}

静态实验建议至少包含三层设置。第一层是 \textbf{单图局部关系评测}：对每个场景采样大量 QRR 和 TRR 问题，分别统计 item accuracy，并按关系类型、对象类别数、遮挡强度和视角复杂度分桶。第二层是 \textbf{多视角整合评测}：向模型提供同一场景的多个观察视角，要求其回答同样的 QRR / TRR 集合，并显式比较单图到多视角的性能变化。第三层是 \textbf{场景级重建评测}：把模型回答转写为距离偏序与方向约束，在统一坐标系中求解一个近似重建布局，再测量其可实现率、冲突率、重建误差、镜像歧义情况以及与真实场景的一致性。

在指标上，单题准确率只能作为最低层统计，真正关键的是 scene-level 指标。建议至少报告：
\begin{itemize}[leftmargin=1.5em]
\item \textbf{Constraint realizability}: 回答集合是否存在统一几何解释。
\item \textbf{Contradiction rate}: 约束之间是否出现环、互斥或跨视角冲突。
\item \textbf{Reconstruction error}: 重建布局与真实场景在相似变换对齐后的结构误差。
\item \textbf{View-consistency gap}: 单图与多视角条件下 scene-level 指标的差值。
\item \textbf{Belief collapse rate}: 多视角输入后从“局部可答”退化为“整体不可实现”的比例。
\end{itemize}

\subsection{当前静态观察}

基于当前已经完成的 540 个场景和 200000+ QA 对，我们已经得到一组对论文中心论点非常关键的初步观察。第一，在单图条件下，QRR 准确率约为 70\%，说明模型对局部相对远近并非完全无感，但这种能力仍明显弱于一个真正稳定的空间系统。第二，当输入扩展为多视角时，QRR 表现并未提升，反而下降到接近随机猜测，这表明模型没有把多个视角整合为更完整的场景表征，而更像是被额外视觉证据破坏了原本脆弱的局部启发式。第三，TRR 长期停留在随机水平，说明模型对统一参照系中的相对方向关系几乎没有形成可靠表征。第四，基于 200K+ QA 对进行微调或强化学习并没有带来显著改善，意味着单纯增加关系监督并不会自动产生几何一致的场景信念。

\subsection{动态场景扩展}

如果论文要从“视觉的根本目的”这一更大叙事切入，那么实验设计不应停留在静态图片。更合理的下一步是引入 \textbf{动态动画场景}，让 benchmark 从静态空间扩展到时空一致性。当前仓库中已有动态图生成与时序关系抽取管线，这为后续实验提供了很好的实现基础，但实验设计本身不需要被现有代码完全限制。

动态场景可以围绕四类任务设计。第一类是 \textbf{时序 QRR}：在不同时间步上询问对象对之间的相对远近，观察模型是否能维持随轨迹变化而更新的偏序结构。第二类是 \textbf{时序 TRR}：询问相对方向关系在运动过程中的演化，测试模型能否维持统一参照系并追踪角度变化。第三类是 \textbf{change-point probing}：直接提问关系何时发生翻转、接近、交叉、相遇、遮挡或离开。第四类是 \textbf{trajectory-consistency probing}：分别根据单帧、短序列、长序列和多视角序列提问，比较模型隐含的场景重建是否在时间推进中仍然相容。

动态实验的核心不应只是“视频题能不能答对”，而应测量模型是否维持了一个 \textbf{随时间更新但结构连续的世界}。因此建议新增以下指标：
\begin{itemize}[leftmargin=1.5em]
\item \textbf{Temporal consistency}: 相邻时间步重建结果的平滑性与因果连续性。
\item \textbf{Relation change accuracy}: 对关系翻转、接近、远离、遮挡事件的识别准确率。
\item \textbf{Future belief stability}: 让模型根据前若干帧预测未来关系，再与真值轨迹对比。
\item \textbf{Cross-time realizability}: 把多个时间步的回答联立后，是否仍存在统一的时空解释。
\end{itemize}

\subsection{干预、对照与能力归因}

为了区分“模型本身具备场景信念”与“模型借助额外流程把题做对”，实验中建议设置两大类 protocol。第一类是 \textbf{原生能力 protocol}：直接回答、简短 CoT、标准多视角输入。第二类是 \textbf{增强能力 protocol}：加入 visual CoT、草图板、程序执行、外部检测器 / 分割器、树搜索或更长的 inference-time scaling \cite{hu2024visualsketchpad,wang2025visuothink,lin2025multimodalcot,suris2023vipergpt,gui2024vpd}。如果增强方法只提升 item accuracy 而不能显著改善 scene-level reconstructability，那么就应将其解释为“解题流程增强”，而不是“真实视觉表征增强”。

\subsection{论文中建议出现的核心表图}

实验部分建议至少包含四类主表。第一类是单图与多视角条件下的 QRR / TRR 准确率表，直接展示从单图到多视角的退化。第二类是 scene-level 指标表，包括 realizability、contradiction rate 与 reconstruction error。第三类是训练干预表，对比 zero-shot、few-shot、SFT、RL 以及 tool-use / visual CoT 方案。第四类是动态实验表，展示 temporal consistency、relation change accuracy 与 cross-time realizability。与之配套的图示建议包括：多视角导致场景信念塌缩的可视化例子、重建成功与失败的场景对比图、以及动态场景中关系随时间翻转的轨迹图。
